\documentclass[10pt,a4paper]{amsart}
\usepackage[margin=19mm]{geometry}
\usepackage{amsmath,amssymb,mathtools}
\usepackage[scr=rsfs,cal=euler]{mathalfa}
\usepackage{xfrac}
\usepackage[unicode,hidelinks,bookmarks=false]{hyperref}
\newcommand{\ie}{i.e.\ }
\newcommand{\cf}{cf.\ }
\newcommand{\resp}{respectively\ }
\newcommand{\Subsection}{\subsection}
\providecommand{\nobreakdash}{\nobreak\hbox{-}}
\setlength{\emergencystretch}{2em}
\multlinegap0pt
\usepackage{bm}
\usepackage{bbm}
\usepackage{dsfont}
\usepackage{xspace}
\usepackage{cancel}
\usepackage{mathtools}
\usepackage[shortlabels]{enumitem}
\usepackage{amsthm}
\makeatletter
\@ifundefined{theorem}{\newtheorem{theorem}{Theorem}}{}
\@ifundefined{conjecture}{\newtheorem{conjecture}[theorem]{Conjecture}}{}
\@ifundefined{lemma}{\newtheorem{lemma}[theorem]{Lemma}}{}
\makeatother
\newcommand{\F}{\mathbb{F}}
\newcommand{\pp}{\mathbb P}
\newcommand{\rad}{\mathrm{rad}}
\title[Minimal Value Set Polynomials]{Minimal Value Set Polynomials}
\author{Herivelto Borges, Lucas Reis}
\date{Source: arXiv:2508.07113v1 (CC BY 4.0)}
\begin{document}
\maketitle

\noindent\emph{Source and licence.} Minimal Value Set Polynomials. Version arXiv:2508.07113v1 (CC BY 4.0), distributed under the Creative Commons Attribution 4.0 International licence, as recorded in the arXiv licence metadata for this exact version and shown on the abstract page https://arxiv.org/abs/2508.07113v1. Full rights notice: licenses/arxiv-2508.07113-CC-BY-4.0.txt.\par\smallskip

\noindent\textit{Editorial scope.} Counted unit: the Theorem whose source label is \texttt{thm:p4-a} in the article (source0.tex lines 816--818, source label \texttt{thm:p4-a}), delivered together with the article's own complete original proof of it (source0.tex lines 820--867, one \texttt{proof} environment of 48 lines). No mathematics is written, repaired or re-derived: statement and proof are the article's own text line for line. Required context retained uncounted: conj (lines 147--155); pre-conj (lines 222--241); abc (lines 305--308); thm:mills (lines 354--362); lem:aux (lines 428--434); lem:gamma (lines 666--668); rad1 (lines 774--776); eq:N0 (lines 783--785); eq:low (lines 804--804). Every definition, display and label of those lines is retained unchanged. Omitted from this package and not redistributed: 1--146, 156--221, 242--304, 309--353, 363--427, 435--665, 669--773, 777--782, 786--803, 805--815, 819--819, 868--1084. Nothing inside a retained excerpt is omitted or altered: every retained body is delivered exactly as the article's lines are written, so no omission receipt of any kind accompanies this package. Citations: the retained excerpts carry no citation command at all, so no bibliography entry of the article is transcribed and no reference is redistributed. Reference adaptation: none was needed and none was made; every reference inside the delivered unit resolves inside it, so no unresolved reference or citation is produced. Printed slips kept literal: no printed word, symbol, hypothesis or number of the counted statement or of its proof was altered. Layout and class: the article's own class is \texttt{amsart} with packages including amsfonts, amssymb, amsmath, amsthm, url, enumerate, xy, color, graphicx, algpseudocode, algorithm, hyperref, todonotes, refcheck; the wrapper is the lane's pinned \texttt{amsart} editorial wrapper at 10pt on A4, which restates the theorem-like environments the retained text needs and adds the article's own macro definitions that the retained text needs (\texttt{\textbackslash F}, \texttt{\textbackslash pp}, \texttt{\textbackslash rad}), with extra wrapper packages (bm, bbm, dsfont, xspace, cancel, mathtools, enumitem). The delivered numbering is the unit's own. Assets: no retained span contains a figure environment, a graphics-inclusion command, a PDF-page inclusion, a table or a TikZ picture; no image file is copied into this package, the package declares no asset dependency and no image file is redistributed with it. Licence: version arXiv:2508.07113v1 of this article is distributed under the Creative Commons Attribution 4.0 International licence, as recorded in the arXiv licence metadata for this exact version and shown on the abstract page https://arxiv.org/abs/2508.07113v1. A full rights notice is delivered at licenses/arxiv-2508.07113-CC-BY-4.0.txt. Characters: every retained excerpt is pure ASCII, so the delivered engine, XeLaTeX with its default Latin Modern text font, reports no missing character.
\begin{conjecture}\label{conj}
Let $q=p^n$, $k$ a divisor of $n$, and $v\geq 1$ a divisor of $p^k-1$.
Let $\mathcal{U} \subseteq \mathbb{F}_q$ be an $\mathbb{F}_{p^k}$-vector space of dimension $m$ such that $1 \in \mathcal{U}$ and
$\# \mathcal{U}^v>2$. If $\mathcal{U}$ is not a field, then $\mathcal{P}\left(\mathcal{U}^v, q\right)=\left\{f^v \mid f \in \mathcal{P}(\mathcal{U}, q)\right\}$. Otherwise,
$$
\mathcal{P}\left(\mathcal{U}^v, q\right)=\left\{\left.f^{\frac{\left(p^e-1\right) v}{p^{m k}-1}} \right\rvert\, f \in \mathcal{P}\left(\mathbb{F}_{p^e}, q\right)\right\},
$$
where $\mathbb{F}_{p^e}  \subseteq \F_q$ is the smallest field containing $\mathcal{U}^v$.
\end{conjecture}

\begin{theorem}\label{pre-conj}  Let $q=p^n$ be a prime power, $k$ a divisor of $n, v>1$  a divisor of $p^k-1$, and let $\mathcal{U} \subseteq \mathbb{F}_q$ be an $\mathbb{F}_{p^k}$-vector space of dimension $m$ such that $\# \mathcal{U}^v>2$ and $1 \in \mathcal{U}$. Then the following hold.
\begin{enumerate}[\rm(i)]
\item If $\mathcal{U} \neq \mathbb{F}_{p^{m k}}$, then $m>1$ and $\operatorname{deg}(F)>\left(p^{n / 2}+1\right) p^{m k}$ for every
$$
F \in \mathcal{P}\left(\mathcal{U}^v, q\right) \backslash\left\{f^v \mid f \in \mathcal{P}(\mathcal{U}, q)\right\} .
$$
\item  If $\mathcal{U}=\mathbb{F}_{p^{m k}}$, let e be the smallest positive integer such that $\mathcal{U}^v \subseteq \mathbb{F}_{p^e} \subseteq$ $\mathbb{F}_{p^{m k}}$, and set $t=\frac{v\left(p^e-1\right)}{p^{m k}-1} $. In this case, if
$$
F \in \mathcal{P}\left(\mathcal{U}^v, q\right) \backslash\left\{f^t \mid f \in \mathcal{P}\left(\mathbb{F}_{p^e}, q\right)\right\},
$$
then $t>1$, $\mathcal{U}^v$ is not a subfield of $\mathbb{F}_q$, and $\operatorname{deg}(F) \geq\left(p^{n / 2}+1\right) p^e$.
 \end{enumerate}
In particular, Conjecture  \ref{conj} holds in each of the following cases:
\begin{enumerate}[\rm(a)]
\item $\# \mathcal{U}^v>p^{n / 2-1}$;
\item  $m>1$ and $(2 m-1) k \geq n / 2$.

 \end{enumerate}

\end{theorem}

\begin{theorem}[Mason-Stothers]\label{abc} Let $k$ be an algebraically closed field and $a, b, c\in k[x]$ pairwise relatively prime polynomials such that $a+b=c$. If the derivatives of these polynomials are not all vanishing, then 
$$\max\{\deg(a), \deg(b), \deg(c)\}\le \deg(\rad(abc))-1,$$
where $\rad(abc)$ denotes the square-free part of $abc$.
\end{theorem}

\begin{theorem}[Mills]\label{thm:mills}
Let $F\in \F_q[x]$ be an MVSP with value set $V_{F}=\{\gamma_0, \ldots, \gamma_r\}$, where $\# V_F=r+1>2$ and $\# (F^{-1}(\gamma_0)\cap \F_q)\le \# (F^{-1}(\gamma_i)\cap \F_q)$ for  $i=1,\ldots, r$. Set $L=\gcd(F-\gamma_0,x^q-x)$  and $T=\prod\limits_{i=0}^{r}(x-\gamma_i)$. Then there exist positive integers $m, k, v$  and polynomials $A, B, N\in \F_q[x]$, where $L\nmid N$  such that the following hold.
\begin{enumerate}[\rm(i)]
\item $\F_{p^k} \subseteq \F_q$, $v$ divides $p^{k}-1$, and $vr+1=p^{mk}$.
\item $F=L^vN^{p^{mk}}+\gamma_0$.
\item $L=A^{p^{mk}}x+B^p$.
\item $T(x+\gamma_0)/x=\sum\limits_{i=0}^{m}w_ix^{\frac{p^{ki}-1}{v}}  \in \F_q[x]$,  with $w_m=1$ and $w_0\ne 0$.
\end{enumerate}
\end{theorem}

\begin{lemma}\label{lem:aux}
Let $\F_{p^k}$ be a subfield of $ \F_q$, $v$ a divisor of $p^k-1$, and $\mathcal U\subseteq \F_q$ an $\F_{p^k}$-vector space with $\#\mathcal U^v>2$. Then the following hold.
\begin{enumerate}[\rm(i)]
\item $F\in\mathcal P(\mathcal U, q)$ if and only if $F^v\in\mathcal P(\mathcal U^v, q)$.
\item If $S=a\cdot\mathcal U^{v}+b$, where $a, b\in \F_q$ and $a\neq 0$, then $\mathcal P(S, q) \neq \emptyset$.
\end{enumerate}
\end{lemma}

 \begin{lemma}\label{lem:gamma}
Let $\F_{p^k}$ be a subfield of $ \F_q$, $v$ a divisor of $p^k-1$, and $\mathcal U\subseteq \F_q$  an $\F_{p^k}$-vector space such that $1\in \mathcal U$ and $\#\mathcal U^v>2$. Suppose that $F\in \mathcal P(\mathcal U^v, q)$, and let $\gamma_0 \in V_F=\mathcal U^v$ be as in Theorem~\ref{thm:mills}. If $v>1$ and $\mathcal U^v$ is not a subfield of $\F_q$, then $\gamma_0=0$.
\end{lemma}

\begin{equation}\label{rad1}
\deg(N)\geq \deg(\rad(N))\ge p^{n/2}+1-2\deg(L) p^{-n/2}.
\end{equation}

\begin{equation}\label{eq:N0}
\deg(N)\ge\frac{r-1}{r+1}p^{n/2}+1.
\end{equation}

Since  $V_F=\left(\F_{p^{mk}}\right)^v$ and $p^e-1=t\cdot \frac{p^{mk}-1}{v}$ divides $q-1=p^n-1$, we have that $F(y)$ is a $t$-th power of an element in $\F_q^*$ for every $y\in \F_q$ with $F(y)\ne 0$. Following the arguments of the previous item, we conclude that $N$ satisfies Eq.~\eqref{eq:N0} and, if $r\ge 3$ or $m_0k_0\le \frac{n}{2}$, then \begin{equation}\label{eq:low}\deg(F)\ge (p^{n/2}+1)p^{m_0k_0}.\end{equation} 

\begin{theorem}\label{thm:p4-a}
 Conjecture~\ref{conj} holds for $q=p^4$. 
\end{theorem}

\begin{proof}
Let $k$ be a divisor of $4$, $v>1$  a divisor of $p^k-1$, and  $\mathcal U\subseteq \F_{p^4}$  an $\F_{p^k}$-vector space of dimension $m$ such that $(m, v)\ne (1, p^k-1)$ and $1\in\mathcal U$. 


If $\mathcal U\neq \F_{p^{mk}}$, then from  (i) and (a) of Theorem~\ref{pre-conj},  Conjecture~\ref{conj} holds for $q=p^4$.  Let us assume  $\mathcal U=\F_{p^{mk}}$, and let $e, t$ be as in assertion (ii) of Theorem~\ref{pre-conj}. Suppose 
\begin{equation}\label{eq-prov}
F\in \mathcal P(\mathcal U^v, p^4)\setminus \{ f^{t}\,|\, f\in \pp(\F_{p^e}, q)\},
\end{equation}
 and let $k_0, v_0,$ and  $m_0$ be the positive integers associated to $F$, as in Theorem~\ref{thm:mills}. Moreover, set $r=\# V_F-1\ge 2$, that is, $r=\frac{p^{mk}-1}{v}$. Observe that any MVSP $F\in \F_{p^4}[x]$ with $\#V_F>1$ satisfies $\deg(F)\le p^4$, and 
recall that $\deg(F)\ge (p^{n/2}+1)p^{m_0k_0}$ was proved in Theorem~\ref{pre-conj} (see Eq. \eqref{eq:low}). Hence $m_0=k_0=1$. We have that $\frac{p^{mk}-1}{v}=\frac{p^{m_0k_0}-1}{v_0}=\frac{p-1}{v_0}$,  and thus $e$ is the least positive integer such that $\frac{p-1}{v_0} $ divides $p^e-1$, that is, $e=1$. In particular, $t=\frac{(p^e-1)v}{p^{mk}-1}=\frac{(p^e-1)v_0}{p^{m_0k_0}-1}=v_0$,  hence $rt+1=p$. Considering \eqref{eq-prov}, assertion (ii) of  Theorem \ref{pre-conj} gives $t\ge 2$, and then  $ r\ge 2$  implies $p\ge 5$. 

From Theorem~\ref{thm:mills}, there exist polynomials $A, B, L, N\in \F_{p^4}[x]$ such that the following hold:
\begin{itemize}%[(i)]
\item $L=\gcd(F,x^q-x)=A^px+B^p$;
\item $F=L^tN^p$;
\item  $rt+1=p$.
\end{itemize}
Note that  we are using that $\gamma_0=0$, which  follows from Lemma~\ref{lem:gamma}, as  (ii) of Theorem~\ref{pre-conj} gives that $\mathcal{U}^v$  is not a subfield of $\mathbb{F}_q$ and  $v_0=t>1$.

Eq. \eqref{eq-prov} and  Lemma~\ref{lem:aux} imply  that $F$, hence $N$, is not of the form $H^t$ with $H\in \F_{p^4}[x]$. 
Observe that  if  every root of $N$ is of multiplicity $\equiv 0, t\pmod p$, then $N=N_1^tN_2^p$, where $N_2$ is not of the form $H^t$. Following the proof of  Eqs. \eqref{rad1} and 
\eqref{eq:N0}   in the proof of Theorem~\ref{pre-conj}, we have that  $N_2$ is not of the form $aH^t$ with $a\in \F_q, H\in \F_q[x]$ and $$\deg(N_2)\ge \max\left\{p^2+1-2\deg(L)p^{-2}, \frac{r-1}{r+1}p^2\right\}.$$
The latter implies 
 $$\deg(F)>p\cdot \deg(N)\ge p^2\deg(N_2)\ge \frac{r-1}{r+1}p^4.$$ 
  Since $F$ is an MVSP with a value set of size $r+1\ge 3$, we have $\deg(F)<\frac{p^4}{2}$, which  contradicts the inequality  $\deg(F)>\frac{r-1}{r+1}p^4$ if $r>2$. If $r=2$, then $t=\frac{p-1}{2}$, since $rt+1=p$. As $p^2\deg(N_2)<\deg(F)<\frac{p^4}{2}$, we conclude that $\deg(N_2)<\frac{p^2}{2}$. Since  $\deg(N_2)\ge p^2-2\deg(L)p^{-2}+1$, we obtain $\deg(L)>\frac{p^4}{4}$. Hence $\deg(F)\ge t\cdot \deg(L)>\frac{p^4(p-1)}{8}\ge \frac{p^4}{2}$, since $p\ge 5$, a contradiction with  $\deg(F)<\frac{p^4}{2}$.  
  
  Recall that $p=rt+1$ with $r, t\ge 2$.  Hence,  once we prove that any root of $N$ has multiplicity $\equiv 0, t\pmod p$, we conclude that Conjecture~\ref{conj} holds for $q=p^4$.  If the derivative $N'$ is vanishing, this is trivially true. Assume that $N'$ is not vanishing.

Under our hypothesis, $V_F=(\mathbb{F}_p)^t=\mathcal C_r\cup\{0\}$, where $\mathcal C_r\subseteq \F_{p}$ is the cyclic group of order $r$. Hence  $T(x)=\prod_{s\in \mathcal V_F}(x-s)=x^{r+1}-x$, and thus $T'(x)=(r+1)x^r-1$. In particular, $T'(a)=r$ for every $a\in V_F\setminus \{0\}$. Since $tr=-1\in \F_p$ and $\# V_F>2$, Theorem~\ref{thm:mills} implies that $$F^{r+1}-F=T(F)=-r(x^{p^4}-x)F'=(x^{p^4}-x)L^{t-1}N^p,$$ and thus replacing $F$ with  $L^tN^p$ yields
\begin{equation}\label{eq:main}\underbrace{(A^{p}\cdot x+B^{p})N^r}_{P_1}=\underbrace{Ax^{p^{3}}+B}_{P_2}.\end{equation}
Observe that for every $b\in \F_{p^4}$, the degree and the value set of $F(x+b)$ coincide with those  of $F$. Moreover, the roots of $N$ have multiplicity $\equiv 0, t\pmod p$ if and only if the same holds for $N(x+b)$. Thus we may assume that $F(0)\ne 0$. Since $L=\gcd(F, x^q-x)=A^px+B^p$, we have $\gcd(A, B)=\gcd(B, x)=1$. In particular, Eq. \eqref{eq:main} implies  \begin{equation}\label{eq:gcd}\gcd(A, N)=\gcd(A, B)=\gcd(B, N)=\gcd(N, x)=\gcd(B, x)=1.\end{equation} If $W(P, Q)=P'Q-PQ'$ denotes the Wronskian map, Eq.\eqref{eq:main} entails that $W(P_1, N^r)=W(P_2, N^r)$, and then
\begin{equation}\label{eq:main2} A^{p}N^{r+1}=S\cdot x^{p^{3}}+T,\end{equation}
where $S=A'N-rAN'$ and $T=B'N-rBN'$. 
Since $N'\neq 0$ and $\gcd(N,A)=\gcd(N,B)=1$, we have that neither $S$ nor $T$ is vanishing. Moreover, since  $\gcd(N,xA)=1$, we have  $$G=\gcd(S x^{p^{3}}, T)=g_1g_2,$$ where $g_1=\gcd(N^{r+1}, T, S)=\gcd(N, N')$ and $g_2=\gcd\left(\frac{Sx^{p^3}}{g_1}, \frac{T}{g_1}, A^p\right)=\gcd\left(\frac{Sx^{p^3}}{g_1}, A^p\right)$, in view of Eq. \eqref{eq:gcd}. Since $g_2$ divides $A^p$, Eq.~\eqref{eq:main2} yields
\begin{equation}\label{eq:main3}\underbrace{\frac{A^{p}}{g_2}\cdot\frac{N^{r+1}}{g_1} }_{R}=\underbrace{\frac{Sx^{p^{3}}}{g_1g_2}}_{S_1}+\underbrace{\frac{T}{g_1g_2}}_{T_1}.\end{equation} 
By construction, the polynomials $R, S_1$, and $T_1$ are pairwise relatively prime. We claim  that $R$ is a polynomial of vanishing derivative.
Our proof relies on Theorem~\ref{abc}.  Therefore, if we set  $$\Delta=\max\{\deg R, \deg S_1, \deg T_1 \}\quad\text{and}\quad \varepsilon=\deg \rad(RS_1T_1),$$
it suffices to prove that $\Delta>\varepsilon-1$. Set $n_1=\deg(N)>0$ and $d=\max\{\deg(A), \deg(B)\}$.
%We have that $(p-1)d+rn_1\le p^3-1$ by Eq.\eqref{eq:main} and $n_1>\frac{1}{3}p^{2}$ by Eq.\eqref{eq:N}. Since $r\ge 2$, the former implies that $d<p^2$. 
Write $A(x)=x^{\ell}\cdot A_0(x)$, where $0\le \ell\le d$, and $A_0\in \F_{p^4}[x]$ is not divisible by $x$. 
Note that if $\ell \geq p^2$, then Eq.~\eqref{eq:main2} implies $T=B'N-rBN'\neq 0$ is divisible by $x^{p^3}$, and then  $\deg B+ \deg N >p^3$, which gives
$$\deg F\geq p(\deg B+\deg N)>p^4,$$
a contradiction. Therefore, $\ell < p^2$. From Eq.~\eqref{eq:gcd}, $N(0)\ne 0$, hence $g_1$ is not divisible by $x$, and thus $\deg(\rad(R))\le d+n_1-\ell+1$. Since  $\ell p<p^3$, it follows that   $g_2(x)$ is of the form $x^{p\ell}\cdot G_2(x)$ with $G_2\in \F_{p^4}[x]$ not divisible by $x$. Therefore, $\deg(\rad(S_1))\le  \deg(S_1)-(p^3-p\ell)+1$ and 
$\deg(T_1)\le \deg(T)-\deg(g_2) \leq d+n_1-1-p\ell$. We conclude that
$$\varepsilon-1\leq  d+n_1-\ell+\deg(S_1)-(p^3-p\ell)+d+n_1-p\ell.$$
Since $\Delta\ge \deg(S_1)$, it suffices to prove that $2d+2n_1< p^3$, which follows from $p\geq 2$ and $p(d+n_1)\leq \deg F<p^4$. Therefore, $\Delta>\varepsilon-1$.
From Theorem~\ref{abc}, $R$ is a $p$-th power. Since $\gcd(A, N)=1$, the same holds for $N_0=\frac{N^{r+1}}{\gcd(N, N')}$, and thus every root of $N_0$ has multiplicity divisible by $p$. Let $\alpha$ be a root of $N$ with multiplicity $j\not\equiv 0\pmod p$. Hence its multiplicity in $N_0$ equals $j(r+1)-(j-1)=jr+1$, forcing that $jr\equiv -1\pmod p$. Since $tr+1=p$, we obtain  $j\equiv t\pmod p$.
\end{proof}

\end{document}
